\chapter*{How to read this manual}
\addcontentsline{toc}{chapter}{How to read this manual}
\markboth{How to read this manual}{}

This document explains the \code{Luna_HiFi} repository as it is used for the
tricycle task: how the code is organised, what every line of every
tricycle-related file does, and how each piece plugs into Isaac Lab, the Newton
physics engine and the \code{rsl_rl} learning library. It is written for
someone who can read Python but has not worked inside Isaac Lab before.
Plain language is used on purpose. Where a word has a precise meaning in Isaac
Lab it is introduced in a box the first time it appears and it is also listed
in the glossary (Appendix~\ref{app:glossary}).

\section*{Two parts}

\textbf{Part~I, the big picture,} explains the system from the outside in: the
repository layout, what reinforcement learning is, what happens between typing
\code{isaaclab train} and a checkpoint appearing on disk, how the physics is
put together, and how the car model becomes a USD asset.

\textbf{Part~II, every file, every line,} walks through each file that takes
part in tricycle training. Each file is shown in slices with its \emph{real
line numbers} (the listings are read straight from the repository at build
time, so what you see is exactly what is committed), and every slice is
followed by a numbered explanation. Lines that belong together (a closing
bracket, a blank line, a run of imports) are explained as a group, but no
line is skipped.

\section*{Conventions}

\begin{itemize}[leftmargin=1.6em]
  \item \code{monospace} marks names that appear literally in code, on the
        command line or in a file path.
  \item ``Line 42'' always refers to the line number printed in the listing
        immediately above it, which is the line number in the actual file.
  \item Coloured boxes carry background knowledge:
\end{itemize}

\begin{isaacbox}{a concept from Isaac Lab}
Blue boxes explain something Isaac Lab does for you, usually with the
file in the Isaac Lab source tree where it is implemented, so you can go and
read the original.
\end{isaacbox}

\begin{newtonbox}{a concept from the physics engine}
Purple boxes explain the Newton physics engine, its MuJoCo-Warp rigid body
solver, or the implicit Material Point Method (MPM) soil solver.
\end{newtonbox}

\begin{whybox}{a design decision}
Green boxes answer ``why is it written like this?''. Most of them record a
lesson that cost real debugging time; they are the same lessons as in
\code{RUNBOOK.md}, section 3, but placed next to the line of code they
explain.
\end{whybox}

\begin{warnbox}{a trap}
Orange boxes flag things that break silently or produce confusing errors.
\end{warnbox}

\section*{What is covered and what is not}

Everything under \srcpath{luna_hifi_tasks/tricycle/}, the package root files,
\code{pyproject.toml}, the converted asset under \srcpath{assets/tricycle/} and
the commands in \code{RUNBOOK.md} are covered line by line. The
\srcpath{luna_hifi_tasks/excavation/} folder is \emph{not} covered. It is a
template for the future rover task, it still contains \code{STUB} and
\code{VERIFY} markers, and nothing in it is imported when the tricycle trains
(Chapter~\ref{ch:repo} shows why).

\section*{Versions this manual matches}

\begin{itemize}[leftmargin=1.6em]
  \item Isaac Lab, \code{develop} branch, early September 2026 (the branch
        that ships the \code{isaaclab train} command, the Newton backend and
        \srcpath{isaaclab_contrib.coupling}). File paths quoted from Isaac Lab
        are relative to the Isaac Lab checkout.
  \item Newton 1.5 (the physics engine Isaac Lab develop installs).
  \item \code{rsl_rl} 5.x as pulled in by \code{isaaclab -i}; the
        configuration in this repository uses the older ``policy'' style and
        Isaac Lab converts it on start-up (Chapter~\ref{ch:agents}).
  \item Python 3.12 on Windows 11, exactly as in \code{RUNBOOK.md}.
\end{itemize}

Quoted line numbers of \emph{Isaac Lab} files may drift as the develop branch
moves; the function and class names are stable enough to search for.
